% GENERATED FILE. Edit canonical lessons, then run npm run generate:books.
% canonical-source-hash: fnv1a64:cc1f258cc76712b6
% canonical-lessons: JA-C328-koukyou, JA-C328-jama, JA-C328-urusai, JA-C328-hikaeru, JA-C328-warikomu, JA-R328-review-323-328

\chapter{Out in Public}
\label{ch:ja-out-in-public}
\hlchaptermodality{328}

\begin{chapteropening}
\textbf{By the end of this chapter:} \emph{I can ask people to behave well in public places and say why.}

Complete the last lesson of chapter 328: I can ask people to behave well in public places and say why.
\end{chapteropening}

\section[\texorpdfstring{\ja{こうきょう} (kōkyō)}{こうきょう (kōkyō)}]{\ja{こうきょう} (kōkyō) — public (places, services)}

\label{lesson:JA-C328-koukyou}



\begin{quote}
Before the new one: say the Japanese for a nuisance, then the Japanese for a bow.
\end{quote}

\subsection*{You'll want to know: \ja{こうきょう}}
\textbf{\ja{こうきょう}} — \emph{kōkyō} — \textquotedblleft{}public (places, services)\textquotedblright{}.

\begin{quote}
\textbf{\ja{こうきょうの} \ja{ばしょは} \ja{みんなが} \ja{つかうから、ごみを} \ja{すてないで} \ja{ください。}} — \emph{kōkyō no basho wa minna ga tsukau kara, gomi o sutenaide kudasai} — \textquotedblleft{}Everyone uses public places, so please do not drop rubbish there.\textquotedblright{}
\end{quote}

\begin{cousinweb}
\textbf{\ja{こうきょう}} (\emph{kōkyō}) is Sino-Japanese (kango): two Chinese characters read with their on'yomi, readings borrowed from Chinese. They are \emph{kō}, \textquotedblleft{}public\textquotedblright{} (Middle Chinese \emph{kuwng}) and \emph{kyō}, \textquotedblleft{}together\textquotedblright{} (\emph{gjowngH}) (KANJIDIC2; Middle Chinese in Baxter's transcription).
\end{cousinweb}

\begin{grammarlens}[title={What you've built}]
One more word for giving reasons and opinions.
\end{grammarlens}

\subsection*{Your turn}
\begin{itemize}
\raggedright
  \item \emph{Say it:} \emph{kōkyō}
  \item \emph{Say it:} \emph{kōkyō}, once more
  \item \emph{Say it:} \emph{kōkyō no basho wa minna ga tsukau kara, gomi o sutenaide kudasai} — Everyone uses public places, so please do not drop rubbish there.
  \item \emph{Recall:} say the Japanese for a nuisance, then the Japanese for a bow, then say \emph{kōkyō} again
\end{itemize}

\begin{tcolorbox}[breakable,colback=teal!4,colframe=teal!35!black,title={Before you move on}]
What does \ja{こうきょう} mean? (\textquotedblleft{}Public (places, services)\textquotedblright{}.) Say it once more.
\end{tcolorbox}
\section[\texorpdfstring{\ja{じゃま} (jama)}{じゃま (jama)}]{\ja{じゃま} (jama) — in the way; a hindrance}

\label{lesson:JA-C328-jama}



\begin{quote}
Before the new one: say the Japanese for a bow, then the Japanese for public.
\end{quote}

\subsection*{You'll want to know: \ja{じゃま}}
\textbf{\ja{じゃま}} — \emph{jama} — \textquotedblleft{}in the way; a hindrance\textquotedblright{}.

\begin{quote}
\textbf{\ja{じてんしゃが} \ja{みちに} \ja{あるから、じゃまです。}} — \emph{jitensha ga michi ni aru kara, jama desu} — \textquotedblleft{}The bicycle is on the road, so it is in the way.\textquotedblright{}
\end{quote}

A na-adjective: \ja{な} before a noun.

\begin{cousinweb}
\textbf{\ja{じゃま}} (\emph{jama}) is Sino-Japanese (kango): two Chinese characters read with their on'yomi, readings borrowed from Chinese. They are \emph{ja}, \textquotedblleft{}wicked\textquotedblright{} (Middle Chinese \emph{zjae}) and \emph{ma}, \textquotedblleft{}demon\textquotedblright{} (KANJIDIC2; Middle Chinese in Baxter's transcription).
\end{cousinweb}

\begin{grammarlens}[title={What you've built}]
One more word for giving reasons and opinions.
\end{grammarlens}

\subsection*{Your turn}
\begin{itemize}
\raggedright
  \item \emph{Say it:} \emph{jama}
  \item \emph{Say it:} \emph{jama}, once more
  \item \emph{Say it:} \emph{jitensha ga michi ni aru kara, jama desu} — The bicycle is on the road, so it is in the way.
  \item \emph{Recall:} say the Japanese for a bow, then the Japanese for public, then say \emph{jama} again
\end{itemize}

\begin{tcolorbox}[breakable,colback=teal!4,colframe=teal!35!black,title={Before you move on}]
What does \ja{じゃま} mean? (\textquotedblleft{}In the way; a hindrance\textquotedblright{}.) Say it once more.
\end{tcolorbox}
\section[\texorpdfstring{\ja{うるさい} (urusai)}{うるさい (urusai)}]{\ja{うるさい} (urusai) — annoying; noisy}

\label{lesson:JA-C328-urusai}



\begin{quote}
Before the new one: say the Japanese for public, then the Japanese for in the way.
\end{quote}

\subsection*{You'll want to know: \ja{うるさい}}
\textbf{\ja{うるさい}} — \emph{urusai} — \textquotedblleft{}annoying; noisy\textquotedblright{}.

\begin{quote}
\textbf{\ja{そとが} \ja{うるさいから、まどは} \ja{あけないで} \ja{ください。}} — \emph{soto ga urusai kara, mado wa akenaide kudasai} — \textquotedblleft{}It is noisy outside, so please do not open the window.\textquotedblright{}
\end{quote}

An i-adjective.

\begin{cousinweb}
\textbf{\ja{うるさい}} (\emph{urusai}) is a native Japanese word (wago), inherited rather than borrowed from Chinese, and its older origin is not settled. JMdict notes that it is usually written in kana (JMdict).
\end{cousinweb}

\begin{grammarlens}[title={What you've built}]
One more word for giving reasons and opinions.
\end{grammarlens}

\subsection*{Your turn}
\begin{itemize}
\raggedright
  \item \emph{Say it:} \emph{urusai}
  \item \emph{Say it:} \emph{urusai}, once more
  \item \emph{Say it:} \emph{soto ga urusai kara, mado wa akenaide kudasai} — It is noisy outside, so please do not open the window.
  \item \emph{Recall:} say the Japanese for public, then the Japanese for in the way, then say \emph{urusai} again
\end{itemize}

\begin{tcolorbox}[breakable,colback=teal!4,colframe=teal!35!black,title={Before you move on}]
What does \ja{うるさい} mean? (\textquotedblleft{}Annoying; noisy\textquotedblright{}.) Say it once more.
\end{tcolorbox}
\section[\texorpdfstring{\ja{ひかえる} (hikaeru)}{ひかえる (hikaeru)}]{\ja{ひかえる} (hikaeru) — to refrain from; to cut down on}

\label{lesson:JA-C328-hikaeru}



\begin{quote}
Before the new one: say the Japanese for in the way, then the Japanese for annoying.
\end{quote}

\subsection*{You'll want to know: \ja{ひかえる}}
\textbf{\ja{ひかえる}} — \emph{hikaeru} — \textquotedblleft{}to refrain from; to cut down on\textquotedblright{}.

\begin{quote}
\textbf{\ja{ねつが} \ja{あるから、おふろは} \ja{ひかえて} \ja{ください。}} — \emph{netsu ga aru kara, ofuro wa hikaete kudasai} — \textquotedblleft{}You have a fever, so please skip the bath.\textquotedblright{}
\end{quote}

Polite \textbf{\ja{ひかえます}}, \emph{hikaemasu}; an ichidan verb. Past (polite): \textbf{\ja{ちちは} \ja{おさけを} \ja{ひかえました。}}, \emph{chichi wa osake o hikaemashita}.

\begin{cousinweb}
\textbf{\ja{ひかえる}} (\emph{hikaeru}) is a native Japanese word (wago), inherited rather than borrowed from Chinese, and its older origin is not settled. Its kanji was attached to it later as a kun'yomi, a native reading given to a Chinese character (JMdict, KANJIDIC2).
\end{cousinweb}

\begin{grammarlens}[title={What you've built}]
One more word for giving reasons and opinions.
\end{grammarlens}

\subsection*{Your turn}
\begin{itemize}
\raggedright
  \item \emph{Say it:} \emph{hikaeru}
  \item \emph{Say it:} \emph{hikaeru}, once more
  \item \emph{Say it:} \emph{netsu ga aru kara, ofuro wa hikaete kudasai} — You have a fever, so please skip the bath.
  \item \emph{Recall:} say the Japanese for in the way, then the Japanese for annoying, then say \emph{hikaeru} again
\end{itemize}

\begin{tcolorbox}[breakable,colback=teal!4,colframe=teal!35!black,title={Before you move on}]
What does \ja{ひかえる} mean? (\textquotedblleft{}To refrain from; to cut down on\textquotedblright{}.) Say it once more.
\end{tcolorbox}
\section[\texorpdfstring{\ja{わりこむ} (warikomu)}{わりこむ (warikomu)}]{\ja{わりこむ} (warikomu) — to cut in line; to push into a queue}

\label{lesson:JA-C328-warikomu}



\begin{quote}
Before the new one: say the Japanese for annoying, then the Japanese for to refrain from.
\end{quote}

\subsection*{You'll want to know: \ja{わりこむ}}
\textbf{\ja{わりこむ}} — \emph{warikomu} — \textquotedblleft{}to cut in line; to push into a queue\textquotedblright{}.

\begin{quote}
\textbf{\ja{れつが} \ja{ながいから、わりこまないで} \ja{ください。}} — \emph{retsu ga nagai kara, warikomanaide kudasai} — \textquotedblleft{}The queue is long, so please do not push in.\textquotedblright{}
\end{quote}

Polite \textbf{\ja{わりこみます}}, \emph{warikomimasu}; a godan verb. Past (polite): \textbf{\ja{おとこが} \ja{わりこみました。}}, \emph{otoko ga warikomimashita}.

\begin{cousinweb}
\textbf{\ja{わりこむ}} (\emph{warikomu}) is a native Japanese word (wago), inherited rather than borrowed from Chinese: it is built from \emph{waru}, \textquotedblleft{}to split\textquotedblright{} and \emph{komu}, \textquotedblleft{}to crowd in\textquotedblright{}, and its kanji were attached later as kun'yomi (JMdict).
\end{cousinweb}

\begin{grammarlens}[title={What you've built}]
One more word for giving reasons and opinions.
\end{grammarlens}

\subsection*{Your turn}
\begin{itemize}
\raggedright
  \item \emph{Say it:} \emph{warikomu}
  \item \emph{Say it:} \emph{warikomu}, once more
  \item \emph{Say it:} \emph{retsu ga nagai kara, warikomanaide kudasai} — The queue is long, so please do not push in.
  \item \emph{Recall:} say the Japanese for annoying, then the Japanese for to refrain from, then say \emph{warikomu} again
\end{itemize}

\begin{tcolorbox}[breakable,colback=teal!4,colframe=teal!35!black,title={Before you move on}]
What does \ja{わりこむ} mean? (\textquotedblleft{}To cut in line; to push into a queue\textquotedblright{}.) Say it once more.
\end{tcolorbox}
\section[(dialogue)]{Looking back: from Saying No Politely to Out in Public}

\label{lesson:JA-R328-review-323-328}



\begin{quote}
Say the Japanese for to hold back and to reject.
\end{quote}

\subsection*{Your turn}
Say these aloud:

\begin{itemize}
\raggedright
  \item to hold back, to reject, impossible, no thank you, not particularly
  \item a grumble, a formal complaint, discontent, to criticize, an excuse
  \item to obey, to go against, to take to heart, after all, freedom
  \item a regulation, to break a rule, a ban, a violation, a punishment
  \item courtesy, behaviour, rude, a nuisance, a bow
  \item public, in the way, annoying, to refrain from, to cut in line
\end{itemize}

\begin{tcolorbox}[breakable,colback=teal!4,colframe=teal!35!black,title={Before you move on}]
To hold back? (\textbf{\ja{えんりょする}}, \emph{enryo suru}.) To reject? (\textbf{\ja{こばむ}}, \emph{kobamu}.) Impossible? (\textbf{\ja{むり}}, \emph{muri}.) No thank you? (\textbf{\ja{けっこう}}, \emph{kekkō}.) Not particularly? (\textbf{\ja{べつに}}, \emph{betsuni}.) A grumble? (\textbf{\ja{もんく}}, \emph{monku}.) A formal complaint? (\textbf{\ja{くじょう}}, \emph{kujō}.) Discontent? (\textbf{\ja{ふへい}}, \emph{fuhei}.) To criticize? (\textbf{\ja{ひはんする}}, \emph{hihan suru}.) An excuse? (\textbf{\ja{いいわけ}}, \emph{iiwake}.) To obey? (\textbf{\ja{したがう}}, \emph{shitagau}.) To go against? (\textbf{\ja{さからう}}, \emph{sakarau}.) To take to heart? (\textbf{\ja{うけとめる}}, \emph{uketomeru}.) After all? (\textbf{\ja{けっきょく}}, \emph{kekkyoku}.) Freedom? (\textbf{\ja{じゆう}}, \emph{jiyū}.) A regulation? (\textbf{\ja{きそく}}, \emph{kisoku}.) To break a rule? (\textbf{\ja{やぶる}}, \emph{yaburu}.) A ban? (\textbf{\ja{きんし}}, \emph{kinshi}.) A violation? (\textbf{\ja{いはん}}, \emph{ihan}.) A punishment? (\textbf{\ja{ばつ}}, \emph{batsu}.) Courtesy? (\textbf{\ja{れいぎ}}, \emph{reigi}.) Behaviour? (\textbf{\ja{ぎょうぎ}}, \emph{gyōgi}.) Rude? (\textbf{\ja{しつれい}}, \emph{shitsurei}.) A nuisance? (\textbf{\ja{めいわく}}, \emph{meiwaku}.) A bow? (\textbf{\ja{おじぎ}}, \emph{ojigi}.) Public? (\textbf{\ja{こうきょう}}, \emph{kōkyō}.) In the way? (\textbf{\ja{じゃま}}, \emph{jama}.) Annoying? (\textbf{\ja{うるさい}}, \emph{urusai}.) To refrain from? (\textbf{\ja{ひかえる}}, \emph{hikaeru}.) To cut in line? (\textbf{\ja{わりこむ}}, \emph{warikomu}.)
\end{tcolorbox}
